\documentclass[12pt,a4paper]{article}

\usepackage[utf8]{inputenc}
\usepackage[T1]{fontenc}
\usepackage[slovak]{babel}
\usepackage{amsmath}
\usepackage{amssymb}
\usepackage{array}
\usepackage{booktabs}
\usepackage{longtable}
\usepackage{geometry}
\usepackage{hyperref}
\usepackage{xcolor}
\usepackage{listings}

\geometry{margin=2.5cm}
\hypersetup{
    colorlinks=true,
    linkcolor=black,
    urlcolor=blue
}

\lstdefinestyle{plain}{
    basicstyle=\ttfamily\small,
    breaklines=true,
    columns=fullflexible,
    keepspaces=true,
    showstringspaces=false,
    frame=single,
    rulecolor=\color{black!35}
}
\lstset{style=plain}

\newcommand{\eps}{\varepsilon}
\newcommand{\First}{\mathrm{FIRST}_1}
\newcommand{\Follow}{\mathrm{FOLLOW}_1}
\newcommand{\code}[1]{\ifmmode\text{\ttfamily #1}\else\texttt{#1}\fi}

\begin{document}

\begin{titlepage}
    \centering
    \vspace*{2cm}
    {\LARGE\bfseries littleXML\par}
    \vspace{0.5cm}
    {\Large Návrh a implementácia tabuľkou riadeného LL(1) analyzátora\par}
    \vfill
    {\large Dokumentácia riešenia\par}
    \vspace{0.5cm}
    {\large \today\par}
\end{titlepage}

\tableofcontents
\newpage

\section{Zadanie a cieľ riešenia}

Cieľom projektu je vytvoriť tabuľkou riadený syntaktický analyzátor zhora nadol pre jazyk \code{littleXML}. Riešenie obsahuje lexikálny analyzátor, LL(1) syntaktický analyzátor, prechodovú tabuľku, výpis krokového logu a jednoduché zotavenie z chýb v režime \code{panic}.

Program pracuje so súbormi. Načíta vstupný súbor, vykoná analýzu a zapíše kompletný report do výstupného súboru. Rovnaký report sa vypíše aj do konzoly.

\subsection{Pôvodná BNF gramatika}

Pôvodná gramatika zo zadania je:

\begin{lstlisting}
xmldokument := [xmldecl] element.
xmldecl := "<?xml" "version =" vernumb "?>".
vernumb := number "." number.
element := emptyelemtag | starttag (words | elements) endtag.
starttag := "<" name ">".
endtag := "</" name ">".
words := {word}.
elements := [element].
emptyelemtag := "<" name "/>".
name := (letter | "_" | ":") {namechar}.
namechar := letter | digit | "." | "-" | "_" | ":".
number := digit {digit}.
word := char {char}.
char := letter | digit | ďalšie znaky okrem "<" a ">".
\end{lstlisting}

Formálne možno gramatiku zapísať ako štvoricu
\[
G = (N, T, P, S),
\]
kde \(N\) je množina neterminálov, \(T\) množina terminálov, \(P\) množina pravidiel a \(S\) štartovací symbol. V našom riešení je štartovací symbol
\[
S = \texttt{xmldokument}.
\]


\section{Transformácia gramatiky}

V tejto časti postupujeme rovnako ako pri ručnej konštrukcii LL(1) analyzátora. Najprv prepíšeme EBNF zápis na obyčajné produkčné pravidlá, potom odstránime spoločné prefixy, vypíšeme medzivýslednú gramatiku a až potom pomocou FIRST/FOLLOW pravých strán pravidiel zistíme konflikty.

\subsection{Prepis EBNF formy na pravidlá}

Zo zadania najprv odstránime EBNF skratky. Voliteľná časť \([X]\) sa prepíše ako alternatíva s \(\eps\), opakovanie \(\{X\}\) sa prepíše rekurziou a terminálne rozsahy písmen a číslic zapíšeme presnými skratkami \code{A..Z}, \code{a..z} a \code{0..9}.

\begin{lstlisting}
xmldokument -> xmldecl element | element
xmldecl -> <?xml version= vernumb ?>
vernumb -> number . number
element -> emptyelemtag | starttag words endtag | starttag elements endtag
starttag -> < name >
endtag -> </ name >
words -> word words | epsilon
elements -> element | epsilon
emptyelemtag -> < name />
name -> namestart name_tail
namestart -> letter | _ | :
name_tail -> namechar name_tail | epsilon
namechar -> letter | digit | . | - | _ | :
letter -> A..Z | a..z
number -> digit number_tail
number_tail -> number | epsilon
digit -> 0..9
word -> char word_tail
word_tail -> word | epsilon
char -> letter | digit | . | - | _ | : | ! | @ | #
\end{lstlisting}

Tento zápis ešte nie je LL(1). Je to iba produkčný tvar pôvodného zadania, z ktorého ďalej odstraňujeme konflikty.

\subsection{Úprava pravidiel}

Po dosadení \code{emptyelemtag} a \code{starttag} do pravidla \code{element} dostaneme tri alternatívy so spoločným začiatkom:

\begin{lstlisting}
element -> < name /> | < name > words endtag | < name > elements endtag
\end{lstlisting}

Všetky pravé strany začínajú symbolom \code{<}, preto by sa do bunky \(M[\code{element}, \code{<}]\) zapisovali tri pravidlá. Spravíme ľavú faktorizáciu spoločného prefixu \code{< name}:

\begin{lstlisting}
element -> < name element_tail
element_tail -> /> | > words endtag | > elements endtag
\end{lstlisting}

V pravidle \code{element\_tail} ešte ostal spoločný prefix \code{>} pri dvoch alternatívach, preto faktorizujeme aj ten:

\begin{lstlisting}
element -> < name element_tail
element_tail -> /> | > element_body
element_body -> words endtag | elements endtag
\end{lstlisting}

Po tejto úprave už vieme rozlíšiť prázdny element podľa \code{/>} a neprázdny element podľa \code{>}. Ešte však musíme overiť, či je jednoznačný aj obsah elementu.

\subsection{Upravené pravidlá pred kontrolou LL(1)}

Po uvedených úpravách dostávame túto medzivýslednú gramatiku. Pravidlá sú očíslované, aby sa dali priamo použiť pri tabuľke FIRST/FOLLOW pravých strán pravidiel.

\begin{lstlisting}
1  xmldokument -> xmldecl element
2  xmldokument -> element
3  xmldecl -> <?xml version= vernumb ?>
4  vernumb -> number . number
5  element -> < name element_tail
6  element_tail -> />
7  element_tail -> > element_body
8  element_body -> words endtag
9  element_body -> elements endtag
10 words -> word words
11 words -> epsilon
12 elements -> element
13 elements -> epsilon
14 endtag -> </ name >
15 name -> namestart name_tail
16 namestart -> letter
17 namestart -> _
18 namestart -> :
19 name_tail -> namechar name_tail
20 name_tail -> epsilon
21 namechar -> letter
22 namechar -> digit
23 namechar -> .
24 namechar -> -
25 namechar -> _
26 namechar -> :
27 letter -> A..Z | a..z
28 number -> digit number_tail
29 number_tail -> number
30 number_tail -> epsilon
31 digit -> 0..9
32 word -> char word_tail
33 word_tail -> word
34 word_tail -> epsilon
35 char -> letter
36 char -> digit
37 char -> .
38 char -> -
39 char -> _
40 char -> :
41 char -> !
42 char -> @
43 char -> #
\end{lstlisting}

\subsection{FIRST/FOLLOW pravých strán pravidiel}

Pri tvorbe LL(1) prechodovej tabuľky sa nepočíta iba FIRST neterminálu, ale FIRST každej pravej strany pravidla. Ak pravá strana obsahuje \(\eps\), použije sa aj FOLLOW ľavej strany pravidla.

Označme si množinu začiatkov textového slova:
\[
W = \{\code{.},\code{-},\code{\_},\code{:},\code{!},\code{@},\code{\#},\code{0..9},\code{A..Z},\code{a..z}\}.
\]

Nasledujúca tabuľka je kompletná tabuľka FIRST/FOLLOW pre pravé strany pravidiel z medzivýslednej gramatiky.

\begin{longtable}{>{\ttfamily}p{0.06\textwidth}p{0.36\textwidth}p{0.30\textwidth}p{0.18\textwidth}}
\toprule
\normalfont Č. & \normalfont Pravá strana & \normalfont FIRST pravej strany & \normalfont FOLLOW pri \(\eps\) / poznámka \\
\midrule
\endhead
1 & xmldecl element & \{\code{<?xml}\} & -- \\
2 & element & \{\code{<}\} & -- \\
3 & <?xml version= vernumb ?> & \{\code{<?xml}\} & -- \\
4 & number . number & \{\code{0..9}\} & -- \\
5 & < name element\_tail & \{\code{<}\} & -- \\
6 & /> & \{\code{/>}\} & -- \\
7 & > element\_body & \{\code{>}\} & -- \\
8 & \textbf{words endtag} & \(W \cup \{\code{</}\}\) & \textbf{konflikt s 9} \\
9 & \textbf{elements endtag} & \{\code{<}, \code{</}\} & \textbf{konflikt s 8} \\
10 & word words & \(W\) & -- \\
11 & epsilon & \{\(\eps\)\} & \(\Follow(\code{words})=\{\code{</}\}\) \\
12 & element & \{\code{<}\} & -- \\
13 & epsilon & \{\(\eps\)\} & \(\Follow(\code{elements})=\{\code{</}\}\) \\
14 & </ name > & \{\code{</}\} & -- \\
15 & namestart name\_tail & \{\code{\_}, \code{:}, \code{A..Z}, \code{a..z}\} & -- \\
16 & letter & \{\code{A..Z}, \code{a..z}\} & -- \\
17 & \_ & \{\code{\_}\} & -- \\
18 & : & \{\code{:}\} & -- \\
19 & namechar name\_tail & \{\code{.}, \code{-}, \code{\_}, \code{:}, \code{0..9}, \code{A..Z}, \code{a..z}\} & -- \\
20 & epsilon & \{\(\eps\)\} & \(\Follow(\code{name\_tail})=\{\code{>},\code{/>}\}\) \\
21 & letter & \{\code{A..Z}, \code{a..z}\} & -- \\
22 & digit & \{\code{0..9}\} & -- \\
23 & . & \{\code{.}\} & -- \\
24 & - & \{\code{-}\} & -- \\
25 & \_ & \{\code{\_}\} & -- \\
26 & : & \{\code{:}\} & -- \\
27 & A..Z | a..z & \{\code{A..Z}, \code{a..z}\} & -- \\
28 & digit number\_tail & \{\code{0..9}\} & -- \\
29 & number & \{\code{0..9}\} & -- \\
30 & epsilon & \{\(\eps\)\} & \(\Follow(\code{number\_tail})=\{\code{?>},\code{.}\}\) \\
31 & 0..9 & \{\code{0..9}\} & -- \\
32 & char word\_tail & \(W\) & -- \\
33 & \textbf{word} & \(W\) & \textbf{FIRST/FOLLOW konflikt s 34} \\
34 & \textbf{epsilon} & \{\(\eps\)\} & \(\Follow(\code{word\_tail})=W\cup\{\code{</}\}\), \textbf{konflikt s 33} \\
35 & letter & \{\code{A..Z}, \code{a..z}\} & -- \\
36 & digit & \{\code{0..9}\} & -- \\
37 & . & \{\code{.}\} & -- \\
38 & - & \{\code{-}\} & -- \\
39 & \_ & \{\code{\_}\} & -- \\
40 & : & \{\code{:}\} & -- \\
41 & ! & \{\code{!}\} & -- \\
42 & @ & \{\code{@}\} & -- \\
43 & \# & \{\code{\#}\} & -- \\
\bottomrule
\end{longtable}

Z tabuľky vidíme hlavný \(\First/\First\) konflikt pri netermináli \code{element\_body}. Pravidlá 8 a 9 obsahujú vo FIRST množine rovnaký terminál \code{</}:

\[
\First(\texttt{words endtag}) \cap \First(\texttt{elements endtag}) = \{\texttt{</}\}.
\]

To znamená, že do jednej bunky prechodovej tabuľky by sa zapisovali dve pravidlá:
\[
M[\texttt{element\_body}, \texttt{</}] = 8
\]
a zároveň
\[
M[\texttt{element\_body}, \texttt{</}] = 9.
\]

Okrem toho je v tabuľke vyznačený aj \(\First/\Follow\) konflikt pri \code{word\_tail}: pravidlo 33 má FIRST množinu \(W\), ale epsilon pravidlo 34 sa zapisuje podľa \(\Follow(\code{word\_tail})\), ktoré tiež obsahuje \(W\):
\[
\First(\texttt{word}) \cap \Follow(\texttt{word\_tail}) = W.
\]
Preto musíme upraviť spôsob, akým reprezentujeme obsah elementu.

\subsection{Odstránenie konfliktov a výsledná gramatika}

Konflikt vznikol preto, že \code{words} aj \code{elements} mohli byť prázdne. Obsah elementu preto rozdelíme priamo podľa toho, akým symbolom začína:

\begin{lstlisting}
element_body -> word endtag | element endtag | endtag
\end{lstlisting}

V implementácii tento tvar zapisujeme kompaktnejšie ako samostatný neterminál \code{element\_content}:

\begin{lstlisting}
element_tail -> /> | > element_content endtag
element_content -> word | element | epsilon
\end{lstlisting}

Tento zápis zodpovedá trom možnostiam: textový obsah, vnorený element alebo prázdny obsah. Zároveň už nepoužívame pôvodné pravidlá \code{words} a \code{elements}, ktoré spôsobovali konflikt.

Výsledná LL(1) gramatika použitá v implementácii je:

\begin{lstlisting}
xmldokument -> xmldecl element | element
xmldecl -> <?xml version= vernumb ?>
vernumb -> number . number
element -> < name element_tail
element_tail -> /> | > element_content endtag
element_content -> word | element | epsilon
endtag -> </ name >
name -> namestart name_tail
namestart -> letter | _ | :
name_tail -> namechar name_tail | epsilon
namechar -> letter | digit | . | - | _ | :
letter -> A..Z | a..z
number -> digit number_tail
number_tail -> number | epsilon
digit -> 0..9
word -> char word_tail
word_tail -> word | epsilon
char -> letter | digit | . | - | _ | : | ! | @ | #
\end{lstlisting}

Použité rozsahy sú skratky pre konkrétne terminály:
\[
\texttt{A..Z} = \{A,\ldots,Z\},
\qquad
\texttt{a..z} = \{a,\ldots,z\},
\qquad
\texttt{0..9} = \{0,\ldots,9\}.
\]
Riešenie nepoužíva skratkové tokeny \code{NAME} a \code{TEXT}. Meno tagu, číslo verzie a textové slovo sa skladajú až syntaktickou analýzou z detailných terminálov. V implementačnej tabuľke sa rozsahy iba technicky mapujú na peek triedy \code{LETTER} a \code{NUMBER}.

\subsection{LL(1) podmienka}

Pre dve alternatívy toho istého neterminálu
\[
A \to \alpha \mid \beta
\]
musí pri LL(1) gramatike platiť
\[
\First(\alpha) \cap \First(\beta) = \emptyset.
\]
Ak niektorá alternatíva vie odvodiť \(\eps\), musí byť zároveň splnená podmienka
\[
\First(\alpha) \cap \Follow(A) = \emptyset
\]
pre epsilon alternatívu. Transformácia gramatiky bola urobená práve preto, aby rozhodovanie podľa jedného symbolu bolo jednoznačné.

\section{FIRST1 a FOLLOW1 množiny}

Výpočet množín \(\First\) a \(\Follow\) vychádza z bežných pravidiel pre LL(1) gramatiky.

Pre terminál \(a \in T\) platí:
\[
\First(a) = \{a\}.
\]

Pre pravidlo
\[
A \to X_1X_2\ldots X_n
\]
sa do \(\First(A)\) pridávajú terminály z \(\First(X_1)\). Ak \(X_1 \Rightarrow^* \eps\), pokračuje sa aj symbolom \(X_2\) atď. Ak celá pravá strana vie odvodiť \(\eps\), potom
\[
\eps \in \First(A).
\]

Pre štartovací symbol platí:
\[
\$ \in \Follow(S).
\]

Pre výpočet FOLLOW množín sa používa pravidlo:
\[
A \to \alpha B \beta
\Rightarrow
\First(\beta) \setminus \{\eps\} \subseteq \Follow(B).
\]
Ak
\[
\beta \Rightarrow^* \eps,
\]
potom platí aj:
\[
\Follow(A) \subseteq \Follow(B).
\]

\subsection{Použité skratky rozsahov}

\begin{align*}
0..9 &= \{0,1,2,3,4,5,6,7,8,9\},\\
A..Z &= \{A,B,C,\ldots,X,Y,Z\},\\
a..z &= \{a,b,c,\ldots,x,y,z\}.
\end{align*}

\subsection{Vypočítané množiny}

\begin{longtable}{>{\ttfamily}p{0.23\textwidth}p{0.33\textwidth}p{0.33\textwidth}}
\toprule
\normalfont Neterminál & \(\First\) & \(\Follow\) \\
\midrule
\endhead
xmldokument & \{\texttt{<?xml}, \texttt{<}\} & \{\texttt{\$}\} \\
xmldecl & \{\texttt{<?xml}\} & \{\texttt{<}\} \\
vernumb & \{\texttt{0},\ldots,\texttt{9}\} & \{\texttt{?>}\} \\
element & \{\texttt{<}\} & \{\texttt{\$}, \texttt{</}\} \\
element\_tail & \{\texttt{/>}, \texttt{>}\} & \{\texttt{\$}, \texttt{</}\} \\
element\_content & \{\(\eps\), \texttt{<}, \texttt{.}, \texttt{-}, \texttt{\_}, \texttt{:}, \texttt{!}, \texttt{@}, \texttt{\#}, \texttt{0..9}, \texttt{A..Z}, \texttt{a..z}\} & \{\texttt{</}\} \\
endtag & \{\texttt{</}\} & \{\texttt{\$}, \texttt{</}\} \\
name & \{\texttt{\_}, \texttt{:}, \texttt{A..Z}, \texttt{a..z}\} & \{\texttt{>}, \texttt{/>}\} \\
namestart & \{\texttt{\_}, \texttt{:}, \texttt{A..Z}, \texttt{a..z}\} & \{\texttt{>}, \texttt{/>}, \texttt{.}, \texttt{-}, \texttt{\_}, \texttt{:}, \texttt{0..9}, \texttt{A..Z}, \texttt{a..z}\} \\
name\_tail & \{\(\eps\), \texttt{.}, \texttt{-}, \texttt{\_}, \texttt{:}, \texttt{0..9}, \texttt{A..Z}, \texttt{a..z}\} & \{\texttt{>}, \texttt{/>}\} \\
namechar & \{\texttt{.}, \texttt{-}, \texttt{\_}, \texttt{:}, \texttt{0..9}, \texttt{A..Z}, \texttt{a..z}\} & \{\texttt{>}, \texttt{/>}, \texttt{.}, \texttt{-}, \texttt{\_}, \texttt{:}, \texttt{0..9}, \texttt{A..Z}, \texttt{a..z}\} \\
letter & \{\texttt{A..Z}, \texttt{a..z}\} & \{\texttt{>}, \texttt{</}, \texttt{/>}, \texttt{.}, \texttt{-}, \texttt{\_}, \texttt{:}, \texttt{!}, \texttt{@}, \texttt{\#}, \texttt{0..9}, \texttt{A..Z}, \texttt{a..z}\} \\
number & \{\texttt{0},\ldots,\texttt{9}\} & \{\texttt{?>}, \texttt{.}\} \\
number\_tail & \{\(\eps\), \texttt{0},\ldots,\texttt{9}\} & \{\texttt{?>}, \texttt{.}\} \\
digit & \{\texttt{0},\ldots,\texttt{9}\} & \{\texttt{?>}, \texttt{>}, \texttt{</}, \texttt{/>}, \texttt{.}, \texttt{-}, \texttt{\_}, \texttt{:}, \texttt{!}, \texttt{@}, \texttt{\#}, \texttt{0..9}, \texttt{A..Z}, \texttt{a..z}\} \\
word & \{\texttt{.}, \texttt{-}, \texttt{\_}, \texttt{:}, \texttt{!}, \texttt{@}, \texttt{\#}, \texttt{0..9}, \texttt{A..Z}, \texttt{a..z}\} & \{\texttt{</}\} \\
word\_tail & \{\(\eps\), \texttt{.}, \texttt{-}, \texttt{\_}, \texttt{:}, \texttt{!}, \texttt{@}, \texttt{\#}, \texttt{0..9}, \texttt{A..Z}, \texttt{a..z}\} & \{\texttt{</}\} \\
char & \{\texttt{.}, \texttt{-}, \texttt{\_}, \texttt{:}, \texttt{!}, \texttt{@}, \texttt{\#}, \texttt{0..9}, \texttt{A..Z}, \texttt{a..z}\} & \{\texttt{</}, \texttt{.}, \texttt{-}, \texttt{\_}, \texttt{:}, \texttt{!}, \texttt{@}, \texttt{\#}, \texttt{0..9}, \texttt{A..Z}, \texttt{a..z}\} \\
\bottomrule
\end{longtable}

\section{Prechodová tabuľka}

Prechodová tabuľka je implementovaná v triede \code{Grammar} a zároveň uložená v súbore \code{table.txt}. Položka tabuľky má tvar:
\[
M[A,a] = A \to \alpha,
\]
kde \(A\) je neterminál na vrchu zásobníka a \(a\) je aktuálny vstupný symbol, v reporte označený ako \code{Peek}.

\begin{longtable}{>{\ttfamily}p{0.24\textwidth}>{\ttfamily}p{0.16\textwidth}>{\ttfamily}p{0.46\textwidth}}
\toprule
\normalfont Neterminál & \normalfont Peek & \normalfont Pravá strana pravidla \\
\midrule
\endhead
xmldokument & <?xml & xmldecl element \\
xmldokument & < & element \\
xmldecl & <?xml & <?xml version= vernumb ?> \\
vernumb & NUMBER & number . number \\
element & < & < name element\_tail \\
element\_tail & /> & /> \\
element\_tail & > & > element\_content endtag \\
element\_content & LETTER & word \\
element\_content & NUMBER & word \\
element\_content & . & word \\
element\_content & - & word \\
element\_content & \_ & word \\
element\_content & : & word \\
element\_content & ! & word \\
element\_content & @ & word \\
element\_content & \# & word \\
element\_content & < & element \\
element\_content & </ & epsilon \\
endtag & </ & </ name > \\
name & LETTER & namestart name\_tail \\
name & \_ & namestart name\_tail \\
name & : & namestart name\_tail \\
namestart & LETTER & letter \\
namestart & \_ & \_ \\
namestart & : & : \\
name\_tail & LETTER & namechar name\_tail \\
name\_tail & NUMBER & namechar name\_tail \\
name\_tail & . & namechar name\_tail \\
name\_tail & - & namechar name\_tail \\
name\_tail & \_ & namechar name\_tail \\
name\_tail & : & namechar name\_tail \\
name\_tail & > & epsilon \\
name\_tail & /> & epsilon \\
namechar & LETTER & letter \\
namechar & NUMBER & digit \\
namechar & . & . \\
namechar & - & - \\
namechar & \_ & \_ \\
namechar & : & : \\
letter & LETTER & LETTER \\
number & NUMBER & digit number\_tail \\
number\_tail & NUMBER & number \\
number\_tail & ?> & epsilon \\
number\_tail & . & epsilon \\
digit & NUMBER & NUMBER \\
word & LETTER & char word\_tail \\
word & NUMBER & char word\_tail \\
word & . & char word\_tail \\
word & - & char word\_tail \\
word & \_ & char word\_tail \\
word & : & char word\_tail \\
word & ! & char word\_tail \\
word & @ & char word\_tail \\
word & \# & char word\_tail \\
word\_tail & LETTER & word \\
word\_tail & NUMBER & word \\
word\_tail & . & word \\
word\_tail & - & word \\
word\_tail & \_ & word \\
word\_tail & : & word \\
word\_tail & ! & word \\
word\_tail & @ & word \\
word\_tail & \# & word \\
word\_tail & </ & epsilon \\
char & LETTER & letter \\
char & NUMBER & digit \\
char & . & . \\
char & - & - \\
char & \_ & \_ \\
char & : & : \\
char & ! & ! \\
char & @ & @ \\
char & \# & \# \\
\bottomrule
\end{longtable}

\section{Lexikálny analyzátor}

Lexikálny analyzátor je implementovaný v triede \code{Tokenizer}. Prechádza vstup zľava doprava a produkuje tokeny, ktoré sú priamo terminálmi gramatiky. Každý token obsahuje typ, hodnotu, riadok, stĺpec a index vo vstupnom texte.

Analyzátor najprv skúša pevné viacznakové a jednoznakové terminály:

\begin{lstlisting}
<?xml  version=  ?>  </  />  <  >  .  -  :  _  !  @  #
\end{lstlisting}

Ak sa nenašiel pevný terminál, vstupný znak sa klasifikuje ako:
\[
\texttt{LETTER}, \quad \text{ak } c \in \{A..Z,a..z\},
\]
alebo
\[
\texttt{NUMBER}, \quad \text{ak } c \in \{0..9\}.
\]

Biele znaky sa preskakujú. To znamená, že vstupy \code{<root/>} a \code{< root />} sú lexikálne rozpoznateľné, hoci syntaktická správnosť sa následne posudzuje podľa tokenovej postupnosti a prechodovej tabuľky.

\subsection{Lexikálne chyby}

Ak sa nájde znak, ktorý nepatrí medzi povolené terminály ani do tried \code{LETTER}/\code{NUMBER}, vytvorí sa diagnostika fázy \code{LEX}. V režime \code{off} analýza končí pri prvej lexikálnej chybe. V režime \code{panic} lexikálny analyzátor preskočí neznámy úsek po najbližší rozpoznateľný token alebo biely znak.

\section{Syntaktický analyzátor}

Syntaktický analyzátor je implementovaný v triede \code{Parser}. Ide o tabuľkou riadený LL(1) analyzátor so zásobníkom.

\subsection{Algoritmus}

Na začiatku je na zásobníku:
\[
[\$, S] = [\texttt{EOF}, \texttt{xmldokument}].
\]

V každej iterácii parser vykoná tieto kroky:

\begin{enumerate}
    \item Zoberie symbol z vrchu zásobníka.
    \item Pozrie aktuálny vstupný token, v reporte označený ako \code{Peek}.
    \item Ak je vrch zásobníka terminál a zhoduje sa s \code{Peek}, token sa spotrebuje.
    \item Ak je vrch zásobníka neterminál, použije sa položka prechodovej tabuľky \(M[A,a]\).
    \item Pravá strana vybraného pravidla sa vloží na zásobník v opačnom poradí.
    \item Ak zásobník aj vstup skončia na \code{EOF} bez diagnostiky, vstup je syntakticky správny.
\end{enumerate}

Jedna iterácia v reporte obsahuje:

\begin{lstlisting}
Zasobnik pred
Vstup pred
Peek
Krok
Zasobnik po
Vstup po
\end{lstlisting}

Položka \code{Peek} obsahuje aj pozíciu tokenu vo vstupnom súbore, napríklad:
\begin{lstlisting}
Peek          : LETTER('a') (riadok 1, stlpec 2)
\end{lstlisting}

\section{Zotavenie z chýb}

Program podporuje dva režimy zotavenia:
\[
\texttt{recovery} \in \{\texttt{off}, \texttt{panic}\}.
\]

\subsection{Režim off}

V režime \code{off} analýza končí pri prvej lexikálnej alebo syntaktickej chybe. Tento režim je vhodný na jednoduché overenie, či je vstup správny.

\subsection{Režim panic}

V režime \code{panic} sa parser pokúsi pokračovať po syntaktickej chybe. Ak je na vrchu zásobníka neterminál \(A\) a pre dvojicu \((A,a)\) neexistuje pravidlo v tabuľke, použije sa synchronizačná množina:
\[
sync(A) = \Follow(A).
\]

Parser preskakuje vstupné tokeny, kým nenájde token zo \(sync(A)\) alebo \code{EOF}. Následne zahodí problémový neterminál zo zásobníka a pokračuje v analýze.

\subsection{Príklad panic recovery}

Vstup:
\begin{lstlisting}
<root><a/><b/></root>
\end{lstlisting}

Gramatika povoľuje najviac jeden vnorený element. Po spracovaní \code{<a/>} parser očakáva neterminál \code{endtag}. Aktuálny vstup je však \code{<b/>}. Pre dvojicu
\[
(\texttt{endtag}, \texttt{<})
\]
neexistuje pravidlo. Použije sa:
\[
\Follow(\texttt{endtag}) = \{\texttt{\$}, \texttt{</}\}.
\]

Parser preto preskočí tokeny zodpovedajúce \code{<b/>} a zastaví sa až na \code{</}. V reporte sa objaví krok podobný:

\begin{lstlisting}
Krok          : panic: skipped <, LETTER, />; sync(endtag)={EOF, </}
\end{lstlisting}

Vstup zostane odmietnutý, ale parser dokáže pokračovať a vypísať viac informácií o priebehu analýzy.

\section{Spustenie programu}

Program sa dá spustiť dvoma spôsobmi. Bežný spôsob je bez argumentov:

\begin{lstlisting}
python .\projekt\main.py
\end{lstlisting}

V tomto režime program vypíše súbory z priečinka \code{projekt/inputs}, používateľ si vyberie vstup a následne zvolí režim zotavenia \code{off} alebo \code{panic}. Report sa uloží do priečinka \code{projekt/outputs}. Názov výstupu sa vytvorí z názvu vstupného súboru a zvoleného režimu.

Program stále podporuje aj priame spustenie so zadaným vstupom a výstupom:

\begin{lstlisting}
python .\projekt\main.py <vstupny_subor> <vystupny_subor> [off|panic]
\end{lstlisting}

Ak sa režim nezadá, použije sa predvolený režim \code{off}.

Príklad bez zotavenia:
\begin{lstlisting}
python .\projekt\main.py .\projekt\inputs\01_valid_empty_element.xml .\projekt\outputs\manual_report.txt off
\end{lstlisting}

Príklad s panic zotavením:
\begin{lstlisting}
python .\projekt\main.py .\projekt\inputs\07_invalid_two_nested_children.xml .\projekt\outputs\panic_report.txt panic
\end{lstlisting}

Výstupný report obsahuje:
\begin{itemize}
    \item vstupný súbor a režim zotavenia,
    \item pôvodný vstup,
    \item zoznam tokenov,
    \item kompletný log syntaktického analyzátora,
    \item problémy validácie,
    \item výsledok \code{SYNTAX OK} alebo \code{SYNTAX ERROR}.
\end{itemize}

\section{Štruktúra implementácie}

\begin{longtable}{>{\ttfamily}p{0.32\textwidth}p{0.58\textwidth}}
\toprule
\normalfont Súbor & \normalfont Význam \\
\midrule
\endhead
main.py & vstupný bod programu, spracovanie argumentov a súborový vstup/výstup \\
grammar.py & terminály, neterminály, FOLLOW množiny a prechodová tabuľka \\
tokenizer.py & lexikálny analyzátor detailných terminálov \\
parser.py & tabuľkou riadený LL(1) syntaktický analyzátor \\
analyzer.py & prepojenie lexikálnej a syntaktickej analýzy \\
token.py & triedy \code{Token} a \code{ErrorLog} \\
table.txt & editovateľná prechodová tabuľka \\
inputs/ & pripravené vstupné súbory na testovanie \\
outputs/ & priečinok pre vygenerované reporty \\
\bottomrule
\end{longtable}

\section{Testovacie vstupy}

\subsection{Pripravené vstupy}

\begin{longtable}{>{\ttfamily}p{0.43\textwidth}p{0.18\textwidth}p{0.29\textwidth}}
\toprule
\normalfont Súbor & \normalfont Výsledok & \normalfont Účel \\
\midrule
\endhead
01\_valid\_empty\_element.xml & OK & prázdny element \\
02\_valid\_xml\_declaration.xml & OK & XML deklarácia a verzia \\
03\_valid\_text\_content.xml & OK & textový obsah elementu \\
04\_valid\_one\_nested\_element.xml & OK & jeden vnorený element \\
05\_valid\_nested\_text.xml & OK & hlbšie vnorenie s textom \\
06\_valid\_mismatched\_names\_by\_grammar.xml & OK & tvar je gramaticky správny, názvy tagov sa neporovnávajú \\
07\_invalid\_two\_nested\_children.xml & ERROR & dva vnorené elementy za sebou \\
08\_invalid\_text\_then\_element.xml & ERROR & miešanie textu a elementu v jednom obsahu \\
09\_invalid\_bad\_version.xml & ERROR & neplatný tvar čísla verzie \\
10\_invalid\_name\_starts\_with\_digit.xml & ERROR & meno tagu začína číslicou \\
11\_invalid\_lexical\_recovery.xml & ERROR & neznámy znak na otestovanie recovery \\
\bottomrule
\end{longtable}

\section{Záver}

Výsledkom je tabuľkou riadený LL(1) analyzátor pre jazyk \code{littleXML}. Gramatika bola transformovaná tak, aby odstránila konflikty spoločného prefixu a umožnila rozhodovanie podľa jedného vstupného symbolu. Implementácia používa detailné terminály namiesto tokenov \code{NAME} a \code{TEXT}, vďaka čomu zodpovedá výpočtom FIRST1/FOLLOW1 nad transformovanou gramatikou. Parser poskytuje kompletný krokový log a podporuje panic recovery založené na FOLLOW množinách.

\end{document}
